\section{Sobolev duality}\label{sec:sobolev}
We use the nonpositive Laplacian
\(\Delta=y^2(\partial_x^2+\partial_y^2)\) and
\(d\mu=dx\,dy/y^2\). Throughout this section \(\Gamma\) denotes the
\emph{effective} image in \(\PSL_2(\R)\) of a finite-index subgroup of
\(\SL_2(\Z)\), or a conjugate thereof. In particular, kernel sums never
count both signs of a matrix. Write \(V=\vol(\Gamma\backslash\HH)\).
For a \(\Gamma\)-invariant discrete set \(\widetilde\Lambda\) with finitely
many orbits, put
\[
 e_z=|\Gamma_z|,\qquad
 \nu_\Lambda=\sum_{z\in\Gamma\backslash\widetilde\Lambda}e_z^{-1}\delta_z,
 \qquad \#\Lambda=\sum_{z\in\Gamma\backslash\widetilde\Lambda}e_z^{-1}.
\]
These weights are necessary at elliptic points.

\begin{lemma}[Resolvent kernel; standard]\label{L:4.1}
There is a nonnegative continuous radial function \(g\) on \(\HH\), with
\[
 g(r)\ll(1+r)e^{-\varphi r},\qquad
 \varphi=\frac{1+\sqrt5}{2},
\]
such that the kernel of \((1-\Delta)^{-2}\) on
\(L^2(\Gamma\backslash\HH)\) is
\(K(z,w)=\sum_{\gamma\in\Gamma}g(\operatorname{dist}(z,\gamma w))\).
The sum converges absolutely.
\end{lemma}
\begin{proof}
% O114-FLAG: The Iwaniec section/theorem numbers here and in the spectral
% expansion are inherited from the source; the book is not in sources/,
% so these locators, unlike DI and Drappeau, have not been independently checked.
The Selberg--Harish-Chandra multiplier is
\(h(t)=(5/4+t^2)^{-2}\); see \cite[\S1.8]{Iw} for the transform and its
inversion, and \cite[Theorem 1.14]{Iw} for automorphisation.
It is even, holomorphic in \(|\Im t|<\sqrt5/2\), and integrable against
\(t\tanh(\pi t)\,dt\). Inversion therefore gives a continuous bounded
point-pair kernel, with
\(g(0)=(4\pi)^{-1}\int_{\R}h(t)t\tanh(\pi t)\,dt\).
For completeness, the decay can also be read directly from the resolvent
kernel. For \(s>1\), the kernel of \((-\Delta+s(s-1))^{-1}\) on \(\HH\) is
\[
 R_s(r)=\frac1{2\pi\sqrt2}\int_r^\infty
       \frac{e^{-(s-1/2)v}}{\sqrt{\cosh v-\cosh r}}\,dv.
\]
Differentiating in \(s(s-1)\), at \(s=\varphi\), gives
\[
 g(r)=\frac1{2\pi\sqrt2(2\varphi-1)}
       \int_r^\infty\frac{v e^{-(\varphi-1/2)v}}
                            {\sqrt{\cosh v-\cosh r}}\,dv.
\]
This proves positivity. For \(r\ge1\), substitute \(v=r+u\) and use
\(\cosh(r+u)-\cosh r\gg e^r(e^u-1)\); the resulting integral is
\(\ll(1+r)e^{-\varphi r}\). At \(r=0\) the numerator cancels the
singularity, giving the same boundedness and continuity as inversion.
An orbit of a discrete cofinite group has \(O_{z,w,\Gamma}(e^R)\) points,
with multiplicity, within distance \(R\). Since \(\varphi>1\), the
periodisation is absolutely convergent. The automorphic kernel theorem
identifies it with the multiplier \(h\), which is precisely
\((1-\Delta)^{-2}\).
\end{proof}

\begin{lemma}[Separated sets]\label{L:4.2}
\status{proved}. If distinct points of \(\widetilde\Lambda\) have distance
at least \(r_0>0\), then
\[
 \|\nu_\Lambda\|_{H^{-2}}^2
 :=\langle\nu_\Lambda,(1-\Delta)^{-2}\nu_\Lambda\rangle
 \le C(r_0)\#\Lambda,
\]
where \(C(r_0)\) is independent of the group and the set.
\end{lemma}
\begin{proof}
Using Lemma~\ref{L:4.1} and cancelling the stabiliser multiplicity in
one variable gives
\[
 \langle\nu_\Lambda,(1-\Delta)^{-2}\nu_\Lambda\rangle
 =\sum_{z\in\Gamma\backslash\widetilde\Lambda}\frac1{e_z}
       \sum_{w\in\widetilde\Lambda}g(\operatorname{dist}(z,w)).
\]
Indeed each lift of \(w\) occurs \(e_w\) times in the kernel sum.
The balls of radius \(r_0/2\) around distinct lifts are disjoint.
Those with centres within distance \(R\) of \(z\) lie in the ball of
radius \(R+r_0/2\); comparison with its area
\(4\pi\sinh^2((R+r_0/2)/2)\) bounds their number by \(O_{r_0}(e^R)\).
Thus each inner sum is at most
\(O_{r_0}(\sum_{m\ge0}(1+m)e^{(1-\varphi)m})=O_{r_0}(1)\).
All rearrangements are justified by positivity and this convergent bound.
\end{proof}

\begin{proposition}[Uniform Sobolev duality]\label{L:4.3}
\status{proved}. Suppose \(P\) is smooth and
\(P,\Delta P,\Delta^2P\in L^2(\Gamma\backslash\HH)\). Put
\(\langle P\rangle=V^{-1}\int P\,d\mu\) and
\(P_0=P-\langle P\rangle\). Under Lemma~\ref{L:4.2}'s separation hypothesis,
\[
 \left|\sum_{z\in\Gamma\backslash\widetilde\Lambda}\frac{P(z)}{e_z}
           -\#\Lambda\langle P\rangle\right|
 \le C(r_0)^{1/2}(\#\Lambda)^{1/2}\|(1-\Delta)P_0\|_2.
\]
\end{proposition}
\begin{proof}
Lemma~\ref{L:4.2} says that \((1-\Delta)^{-1}\nu_\Lambda\in L^2\)
with squared norm at most \(C(r_0)\#\Lambda\). Self-adjointness, first
for smooth regularisations of the point masses and then by Sobolev
duality, gives
\[
 \nu_\Lambda(P_0)=
 \langle(1-\Delta)P_0,(1-\Delta)^{-1}\nu_\Lambda\rangle.
\]
Point evaluation agrees with this distributional pairing by local
Sobolev embedding in dimension two. Cauchy--Schwarz proves the claim.
\end{proof}

\begin{corollary}[Heegner forms]\label{L:4.4}
\status{proved}. Let \(\cF\subset\cF_d\) be invariant under a subgroup
\(\Gamma\subset\Gamma_0(d)\) acting on \(w_Q=z_Q/d\), with finitely many
orbits. Proposition~\ref{L:4.3} holds with an absolute constant for its
root set. If the effective stabiliser of infinity consists of translations
by \(h\Z\), and \(\psi\in C_c^\infty(\HH)\), set
\[
 \psi_{\rm per}(w)=\sum_{m\in\Z}\psi(w+mh),\qquad
 P_\psi(w)=\sum_{\gamma\in\Gamma_\infty\backslash\Gamma}
                           \psi_{\rm per}(\gamma w).
\]
Then
\[
 \sum_{z\in\Gamma\backslash\widetilde\Lambda}\frac{P_\psi(z)}{e_z}
  =\sum_{Q\in\cF}\psi(w_Q),\qquad
 \langle P_\psi\rangle=V^{-1}\int_{\HH}\psi\,d\mu.
\]
The same statements hold after conjugating the group by a real dilation.
\end{corollary}
\begin{proof}
The dilation is a hyperbolic isometry, so Lemma~\ref{L:2.2} supplies the
absolute separation \(r_0=\operatorname{arccosh}(3/2)\).
Unfold the group sum, cancelling elliptic stabilisers as in
Lemma~\ref{L:4.2}, and then unfold translations. The root map is
injective at fixed discriminant, so this counts forms, not just roots.
Unfolding the integral proves the mean formula. The Poincar\'e series
and its Laplacian iterates are smooth and vanish sufficiently high in
every cusp: for a nonstabilising coset the transformed height tends to
zero, below the support of \(\psi\), whereas at infinity the stabilising
coset eventually lies above its support. Thus they satisfy the
integrability assumptions of Proposition~\ref{L:4.3}.
\end{proof}

\section{The spectral large sieve and the variance of box Poincar\'e series}
\label{sec:sieve}
We recall the coordinates of Section~\ref{sec:heegner}, keeping the parity
condition throughout. For odd squarefree \(q\) with \((q,d)=1\), conjugation
by \(u=w/(2q)\) takes the effective image of
\(\Gamma'_q=\Gamma_0(2d)\cap\Gamma(2q)\) to
\[
 \Gamma_{M,q}=\left\{\pm\begin{pmatrix}p&t\\r&s\end{pmatrix}
       \in\Gamma_0(M):p\equiv s\equiv1\pmod{2q}\right\},
 \qquad M=4dq^2.
\]
Its cusp at infinity has width one. Set
\(h_q=[\Gamma_0(M):\Gamma_{M,q}]\), with effective groups understood:
\(h_1=1\) and \(h_q=\varphi(q)/2\) for \(q>1\).
The quotient is \((\Z/q\Z)^\times/\{\pm1\}\); hence its characters are
exactly the even characters modulo \(q\), equivalently modulo \(2q\).

\subsection*{The imported spectral estimates}
We specify normalisations to distinguish a cusp width from the parameter
in the large sieve. At a singular cusp \(\mathfrak a\), choose a scaling
matrix giving period one, and write the expansion of an orthonormal
weight-zero Maass form as
\begin{equation}\label{B:eq:fourier}
 u_j(\sigma_{\mathfrak a}z)
 =\sqrt y\sum_{n\ne0}\rho_{j\mathfrak a}(n)
                 K_{it_j}(2\pi|n|y)e^{2\pi inx}.
\end{equation}
This is DI's normalisation \cite[(1.34), p.~231]{DI}.
The Eisenstein coefficients \(\rho_{\mathfrak c\mathfrak a}(n,t)\)
will use the identical \(\sqrt yK_{it}\) convention. The DI cusp
parameter for \(\Gamma_0(M)\) is
\[
 \mu(\mathfrak a)=\frac{(v,M/v)}M
 \quad\text{if }\mathfrak a\sim u/v,\ v\mid M,\ (u,v)=1;
 \qquad \mu(\infty)=M^{-1}.
\]
It is \emph{not} the literal translation width at infinity, which is one.
For \(K\ge1\), \(H\ge1/2\), and \(a_n\) supported on \(H<n\le2H\),
DI \cite[Theorem 2, (1.29)--(1.30), p.~230]{DI} bounds the cuspidal
expression
\[
 \sum_{|t_j|\le K}\frac1{\cosh(\pi t_j)}
       \left|\sum_n a_n\rho_{j\mathfrak a}(n)\right|^2
\]
and its continuous-spectrum analogue by
\((K^2+\mu(\mathfrak a)H^{1+\varepsilon})\sum_n|a_n|^2\), for trivial
nebentypus. With an even character of conductor \(q_0\mid q\), the bound is
\begin{equation}\label{B:eq:large-sieve}
 \ll_\varepsilon
 (K^2+q_0^{1/2}\mu(\mathfrak a)H^{1+\varepsilon})\sum_n|a_n|^2.
\end{equation}
Here the continuous analogue means
\(\sum_{\mathfrak c}(4\pi)^{-1}\int_{-K}^K
\cosh(\pi t)^{-1}|\sum_n a_n\rho_{\mathfrak c\mathfrak a}(n,t)|^2dt\),
summed over all singular cusps. This is
\cite[Proposition 4.7, (4.24)--(4.25)]{Dr}, in the arXiv numbering.
Drappeau uses Whittaker coefficients: since
\(W_{0,it}(4\pi|n|y)=2\sqrt{|n|y}K_{it}(2\pi|n|y)\),
\(\rho^{\rm DI}(n)=2\sqrt{|n|}\rho^{\rm Dr}(n)\).
Thus his \(\sqrt n\rho^{\rm Dr}(\pm n)\) is exactly our coefficient
up to the constant 2; his factor \((1+|t|)^{\pm\kappa}\) is one because
\(\kappa=0\). These estimates hold separately for either sign of \(n\).

For later comparison, DI \cite[Theorem 5, (1.38), p.~232]{DI} states,
for trivial character, \(X\ge1\), and exceptional parameters
\(t_j=i\sigma_j\), \(0<\sigma_j<1/2\),
\begin{equation}\label{B:eq:exceptional-sieve}
 \sum_{j\ {
m exc}}X^{2\sigma_j}
   \left|\sum_{H<n\le2H}a_n\rho_{j\mathfrak a}(n)\right|^2
 \ll_\varepsilon
 (1+\sqrt{\mu(\mathfrak a)HX})
 (1+\sqrt{\mu(\mathfrak a)H^{1+\varepsilon}})\|a\|_2^2.
\end{equation}
The nebentypus analogue is \cite[Lemma 4.8]{Dr}, with second factor
\(1+(q_0\mu(\mathfrak a)H)^{1/2+\varepsilon}\).
The displayed DI statements have been checked against the scan; the
Drappeau numbering and coefficient conversion refer to arXiv:1504.05549,
not to an unchecked published numbering.

\subsection*{The box variance}
Take \(\phi\in C_c^\infty([-2,-1])\), \(W\in C_c^\infty([1,2])\), and
\[
 \psi(u)=\phi(x/\lambda)W(y/Y),\qquad
 L_* =\log(2+(\lambda Y)^{-1}+Y^{-1}+M).
\]
Assume \(\lambda>0\) and \(Y\le\lambda L_*^{-3}\). Periodisation is always
understood; no restriction \(\lambda<1\) is imposed. For nonnegative
nonzero weights its area is \(\asymp\lambda/Y\). In applications
\(\lambda\asymp A/(qF)\), \(Y\asymp1/(qF\sqrt d)\), and \(L_*\ll\cL\).

\begin{proposition}[Box variance]\label{L:5.1}
\status{conditional} on (SEL). On \(\Gamma_{M,q}\backslash\HH\),
\begin{equation}\label{B:eq:variance}
 \|(1-\Delta)(P_\psi-\langle P_\psi\rangle)\|_2^2
 \ll_{\phi,W,\varepsilon}L_*^C
 \left\{\frac\lambda Y
   \left(1+\frac{q^{1/2}}M(L_*/\lambda)^{1+\varepsilon}\right)
       +\frac{\lambda^2}Y\right\}.
\end{equation}
The constant is uniform in \(d,q,\lambda,Y\), and depends on only finitely
many seminorms of the two fixed weights.
\end{proposition}
\begin{proof}
\emph{Reduction and spectral expansion.}
Differentiation commutes with periodisation and
\[
 (1-\Delta)\psi
 =\phi(x/\lambda)(W(v)-v^2W''(v))
       -(Y/\lambda)^2\phi''(x/\lambda)v^2W(v),\qquad v=y/Y.
\]
Moreover \(\int\Delta\psi\,d\mu=0\). Since \(Y/\lambda\le1\), it suffices
to bound \(\|P_\psi-\langle P_\psi\rangle\|_2^2\) by the right side of
\eqref{B:eq:variance}, for arbitrary fixed smooth weights of this type.
For each even \(\chi\bmod q\), let
\[
 P_\chi=\sum_{\gamma\in\Gamma_\infty\backslash\Gamma_0(M)}
            \overline{\chi(\gamma)}\psi_{\rm per}(\gamma u).
\]
Character orthogonality gives \(P_\psi=h_q^{-1}\sum_\chi P_\chi\) and
\begin{equation}\label{B:eq:char-norm}
 \|P_\psi-\langle P_\psi\rangle\|_{\Gamma_{M,q}}^2
 =\frac1{h_q}\sum_\chi
       \|P_\chi-\ind_{\chi=1}\langle P_1\rangle\|_{\Gamma_0(M),\chi}^2.
\end{equation}
The spectral expansion, with Eisenstein measure \(dt/(4\pi)\), is the
usual Parseval formula \cite[Theorem 7.3]{Iw} (with the character
spaces and singular cusps as in \cite[\S4.1]{Dr}).
\emph{This is the sole use of (SEL):} after removing the constant,
all discrete parameters are real, and there is no nonconstant residual
spectrum. Thus the squared norm for a fixed character is the sum of
\(|\langle P_\chi,u_j\rangle|^2\) over cusp forms and
\((4\pi)^{-1}\sum_{\mathfrak c}\int_\R
|\langle P_\chi,E_{\mathfrak c}(\cdot,1/2+it)\rangle|^2dt\).

\emph{Unfolding and Mellin inversion.}
Write \(\widehat\phi(\xi)=\int_\R\phi(v)e^{-2\pi i\xi v}dv\).
The periodised box has Fourier coefficient
\(\lambda\widehat\phi(\lambda n)W(y/Y)\), even when \(\lambda>1\).
Using \eqref{B:eq:fourier}, unfolding gives
\[
 \langle P_\chi,u_j\rangle
 =\sum_{n\ne0}\overline{\rho_j(n)}\lambda\widehat\phi(\lambda n)I_{t_j}(n),
 \quad I_t(n)=\int_0^\infty W(y/Y)y^{-3/2}K_{it}(2\pi|n|y)\,dy.
\]
Set
\[
 G_t(s)=\Gamma((s+it)/2)\Gamma((s-it)/2),\qquad
 \mathcal W(s)=\int_0^\infty W(v)v^{-3/2-s}dv.
\]
Write also
\[
 B_j(s)=\sum_{n\ne0}\overline{\rho_j(n)}
              \lambda\widehat\phi(\lambda n)|n|^{-s}.
\]
The function \(\mathcal W\) is entire and rapidly decreasing vertically,
uniformly on fixed strips. Mellin inversion of the Bessel function gives,
for \(\sigma>0\),
\begin{equation}\label{B:eq:mellin}
 \langle P_\chi,u_j\rangle
 =\frac{Y^{-1/2}}{8\pi i}\int_{(\sigma)}
        G_{t_j}(s)\mathcal W(s)(\pi Y)^{-s}B_j(s)\,ds.
\end{equation}
Absolute convergence follows from the rapid decay of
\(\widehat\phi\), the polynomial growth of Fourier coefficients, and
Stirling's formula.

\emph{Uniform spectral bounds.}
Put \(\sigma_0=1/L_*\). Stirling's formula, including the simple poles
of \(\Gamma\), gives
\begin{align*}
 |G_t(\sigma_0+iv)|&\ll L_*^2e^{-\pi|v|/2} &&(|t|\le1),\\
 |G_t(-1+iv)|&\ll e^{-\pi|t|/2}(1+|t|)^{-2}(1+|v|)^4
                                      &&(|t|\ge1).
\end{align*}
For the second bound, if \(|v|\le|t|/2\) both gamma arguments have
imaginary part comparable to \(|t|\); otherwise
\(1+|t|\ll1+|v|\) and the exponential is
\(e^{-\pi\max(|v|,|t|)/2}\). On this line both real parts are \(-1/2\),
a fixed distance from gamma poles. Since \(Y^{-\sigma_0}\ll1\),
Cauchy--Schwarz in \eqref{B:eq:mellin} yields, for \(|t_j|\le1\),
\[
 |\langle P_\chi,u_j\rangle|^2
 \ll Y^{-1}L_*^4\int_\R |\mathcal W(\sigma_0+iv)|
                                    |B_j(\sigma_0+iv)|^2\,dv.
\]
For \(|t_j|\ge1\), shift to \(\Re s=-1\), crossing only \(s=\pm it_j\).
The residue contribution is
\[
 \frac{Y^{-1/2}}2\sum_\pm
       \Gamma(\pm it_j)\mathcal W(\pm it_j)(\pi Y)^{\mp it_j}B_j(\pm it_j).
\]
Let \(\widetilde B_j(s)=(\pi Y)^{-s}B_j(s)\).
Using \(|\Gamma(it)|^2=\pi/(|t|\sinh(\pi|t|))\), for any fixed large
\(B\) the squared coefficient is at most
\begin{align}\label{B:eq:high-t}
 \frac{C_BY^{-1}}{\cosh(\pi t_j)}\bigg\{&
 (1+|t_j|)^{-B}\sum_\pm|B_j(\pm it_j)|^2\notag\\
 &+(1+|t_j|)^{-4}\int_\R
  |\mathcal W(-1+iv)|(1+|v|)^4|\widetilde B_j(-1+iv)|^2\,dv\bigg\}.
\end{align}

\emph{The large sieve, including both tails.}
Split \(n\) by sign and into dyadic blocks \(H<|n|\le2H\),
\(H=2^{i-1}\), \(i\ge0\). With
\(w_H=1+|\log(\lambda H)|\), one has \(\sum_Hw_H^{-2}\ll1\).
Cauchy--Schwarz across blocks and \eqref{B:eq:large-sieve} imply
\[
 \mathcal S(K;b):=\sum_{|t_j|\le K}\frac{|\sum_{n\ne0}\overline{\rho_j(n)}b_n|^2}
                                             {\cosh(\pi t_j)}
 \ll_\varepsilon\sum_H w_H^2
       (K^2+q^{1/2}M^{-1}H^{1+\varepsilon})\|b^{(H)}\|_2^2.
\]
Complex conjugation of the coefficient vector causes no change in this
inequality; either sign is covered by Drappeau's statement (or reflection
and replacement of \(\chi\) by \(\bar\chi\)). For
\(b_n(s)=\lambda\widehat\phi(\lambda n)|n|^{-s}\), at
\(\Re s=\sigma_0\) or \(0\), the block norm is
\(\ll_B\lambda^2H(1+\lambda H)^{-B}\).
Summing the geometric tails gives, with
\(Z=q^{1/2}M^{-1}\lambda^{-1-\varepsilon}\),
\begin{equation}\label{B:eq:block-sieve}
 \mathcal S(K;b(s))\ll L_*^2\lambda(K^2+Z).
\end{equation}
At \(\Re s=-1\), the coefficients of \(\widetilde B\) have modulus
\(\pi Y|n|\lambda|\widehat\phi(\lambda n)|\); their bound is
\eqref{B:eq:block-sieve} multiplied by \((Y/\lambda)^2\le1\), after
increasing the decay exponent. This also proves the claims when
\(\lambda>1\): then every nonzero mode lies in the rapidly decreasing tail.

For the varying vectors \(b_n(\pm it_j)\), on each unit interval
\([k,k+1]\) use the elementary inequality
\[
 |B_j(\pm it)|^2\le2|B_j(\pm ik)|^2+
              2\int_k^{k+1}|\partial_\tau B_j(\pm i\tau)|^2d\tau.
\]
Differentiating multiplies a coefficient by \(\log|n|\), bounded on a
block by \(O(L_*+|\log(\lambda H)|)\); the extra factors are absorbed
in the same block sums. Consequently the sum on \(K<|t_j|\le2K\)
for these vectors is \(\ll K L_*^4\lambda(K^2+Z)\).
Insert these bounds in \eqref{B:eq:high-t} and the small-\(t\) estimate.
Dyadic summation over \(K\ge1\) converges, since the integral term has
weight \(K^{-4}\) and the residue term has arbitrary weight \(K^{-B}\).
Vertical integrals converge by the decay of \(\mathcal W\). We obtain
\[
 \sum_j|\langle P_\chi,u_j\rangle|^2
       \ll L_*^C\frac\lambda Y(1+Z).
\]
No truncation at logarithmic frequency or height was made.

\emph{The continuous spectrum and its zero mode.}
For nonzero modes the identical argument uses the continuous version of
\eqref{B:eq:large-sieve}. In the unit-interval argument freeze the
Eisenstein parameter in the Fourier coefficients and differentiate only
the auxiliary variable in \(|n|^{-i\tau}\), then evaluate on the diagonal
\(\tau=t\); no derivative of an Eisenstein coefficient is required.
The constant term at infinity is
\(\delta_{\mathfrak c\infty}y^{1/2+it}
 +\Phi_{\mathfrak c\infty}(1/2+it)y^{1/2-it}\).
Its pairing with the box is
\[
 \lambda\widehat\phi(0)Y^{-1/2}
 \{\delta_{\mathfrak c\infty}Y^{-it}\mathcal W(it)
  +\overline{\Phi_{\mathfrak c\infty}(1/2+it)}Y^{it}\mathcal W(-it)\}.
\]
Unitarity of the scattering matrix on the unitary axis gives
\(\sum_{\mathfrak c}|\Phi_{\mathfrak c\infty}(1/2+it)|^2=1\).
Squaring, integrating and summing therefore costs \(O(\lambda^2/Y)\),
independently of the number of cusps. The cross term with nonzero modes
is absorbed by \(|a+b|^2\le2|a|^2+2|b|^2\).
Finally \eqref{B:eq:char-norm} averages exactly \(h_q\) character bounds;
there is no loss from their number. Enlarging
\(\lambda^{-1-\varepsilon}\) to \((L_*/\lambda)^{1+\varepsilon}\) proves
\eqref{B:eq:variance}.
\end{proof}

Without (SEL), the terms with \(t_j=i\sigma_j\) remain. On a block
\(|n|\asymp H\), the Bessel factor has additional size
\((HY)^{-\sigma_j}\). Thus \eqref{B:eq:exceptional-sieve} is used with
\(X=(HY)^{-1}\) when \(HY\le1\), not with \(X=Y^{-1}\).
It does not recover Proposition~\ref{L:5.1} in the required cells;
Section~\ref{sec:uncond} records precisely the resulting unconditional
limitation. Neither the Sobolev argument nor the regular-spectrum large
sieve depends on (SEL).

\section{Per-\texorpdfstring{$d$}{d} counts}\label{sec:perd}
Fix \(c\ge1\). The Type~I subset is
\[
 \cF_d^I=\{[f,4ad,de]:ef-4a^2d=1,\ e,f>0,\ a\in\Z\},
 \qquad n(Q)=cB-A=4acd-f.
\]
Here \(A\) in \([A,B,C]\) is a form coefficient, not the dyadic parameter.
This subset is invariant under \(\Gamma'_1=\Gamma_0(2d)\cap\Gamma(2)\)
on the \(w\)-side. Indeed, the corresponding matrix on the \(z\)-side
has upper-right entry divisible by \(2d\) and lower-left entry even, so
\(B'=2Apt+B(ps+tr)+2Crs\equiv0\pmod{4d}\).
Reduction modulo \(q\) shows that the additional condition \(q\mid n(Q)\)
is invariant under \(\Gamma'_q\). These are the groups of
Section~\ref{sec:sieve}; using \(\Gamma_0(d)\) alone would lose parity.

For a prime \(\ell\nmid2d\), put \(\chi_d(\ell)=(-d/\ell)\) and define
\(g_{c,d}(\ell)\) as the proportion on the quadric
\(B^2-4AC=-4d\) satisfying \(cB-A=0\). Directly,
\begin{equation}\label{B:eq:density}
 g_{c,d}(\ell)=
 \begin{cases}
 (\ell-1)/(\ell^2+\chi_d(\ell)\ell),&\ell\nmid c,\\
 (1+\chi_d(\ell))/(\ell+\chi_d(\ell)),&\ell\mid c.
 \end{cases}
\end{equation}
Extend multiplicatively to squarefree \(q\) coprime to \(2d\).
The denominator counts the quadric: for \(A\ne0\) there are
\(\ell(\ell-1)\) choices, while \(A=0\) gives
\(\ell(1+\chi_d(\ell))\). If \(\ell\nmid c\), the imposed equation
forces \(B\ne0\), giving exactly \(\ell-1\) choices. If \(\ell\mid c\),
it forces \(A=0\), giving the second numerator. Primes dividing \(2d\)
never divide \(n\): both \(f\) and \(n\) are odd, and \(f\) is a unit
at every prime dividing \(d\).

\begin{lemma}[Factorisation of the orbital main term]\label{L:6.1}
\status{proved}. Let \(\#\Lambda_d(q)\) be the orbifold count of
\(\{Q\in\cF_d^I:q\mid n(Q)\}\) modulo \(\Gamma'_q\), and put
\(V(q)=\vol(\Gamma'_q\backslash\HH)\). For squarefree \((q,2d)=1\),
\[
 \frac{\#\Lambda_d(q)}{V(q)}
 =g_{c,d}(q)\frac{\#\Lambda_d(1)}{V(1)}.
\]
\end{lemma}
\begin{proof}
There is nothing to prove for \(q=1\). Reduction of \(\Gamma'_1\)
onto \(\SL_2(\Z/q\Z)\) is surjective: it contains \(\Gamma(2d)\),
and the assertion follows from CRT and surjectivity of reduction for
\(\SL_2(\Z)\). The kernel is \(\Gamma'_q\). Thus the index of effective
groups is
\(I_q=|\SL_2(\Z/q\Z)|/2\), since \(-I\) belongs to \(\Gamma'_1\)
but not to \(\Gamma'_q\).
At a prime \(\ell\mid q\), the stabiliser of a nondegenerate binary form
in \(\SL_2(\mathbb F_\ell)\) is its special orthogonal torus, of order
\(\ell-\chi_d(\ell)\): it is either the split torus or the norm-one
subgroup of \(\mathbb F_{\ell^2}^{\times}\).
Its orbit consequently has \(\ell^2+\chi_d(\ell)\ell\) elements, the
size of the entire quadric. The fraction satisfying the sieve condition
is therefore \(g_{c,d}(\ell)\), independently of the original orbit.
CRT gives the product of these fractions.
For an orbit with stabiliser order \(e_Q\), the weighted number of
sieved \(\Gamma'_q\)-orbits above it is
\(e_Q^{-1}I_qg_{c,d}(q)\), by orbit--stabiliser (equivalently, sum over
the cosets before dividing by the stabiliser). Summing over the original
orbits, and using \(V(q)=I_qV(1)\), proves the formula.
\end{proof}

Choose a fixed physical box, so that in the coordinates \(u=w/(2q)\)
its weights are \(\psi_q(u)=\psi_1(qu)\). Accordingly
\(\lambda_q=\lambda_1/q\), \(Y_q=Y_1/q\), and
\(\int\psi_q\,d\mu=\int\psi_1\,d\mu\). Define
\[
 \mathfrak M_d=\frac{\#\Lambda_d(1)}{V(1)}\int_{\HH}\psi_q\,d\mu,
 \qquad \kappa_c(q)=2^{\omega((q,c))}.
\]
The first quantity is independent of the sieve modulus. Here and below
\(\omega(m)\) counts the distinct prime factors of \(m\).

% O114-FLAG: The source's Theorem 6.2 drops a factor at primes dividing c.
% Its assertion #Lambda_d(q) << q^2 #Lambda_d(1), uniformly in c,q,
% does not follow for (q,c)>1: a split prime dividing c contributes
% |SL_2(F_ell)| g(ell) = 2 ell(ell-1). The theorem below retains
% kappa_c(q)^{1/2}; the source's bound is recovered for (q,c)=1,
% which is the sieve range excluding primes dividing 2cd in the assembly.
\begin{theorem}[Per-\(d\) count]\label{L:6.2}
\status{conditional} on (SEL). Let \(q\) be squarefree, \((q,2d)=1\),
and let \(\psi_q\) be a box as in Proposition~\ref{L:5.1}, with
\(\lambda_q\asymp A/(qF)\), \(Y_q\asymp1/(qF\sqrt d)\), \(A,F\ge1\).
Assume the parameters are of polynomial size in \(N\), so \(L_*\ll\cL\).
Then, for \(0<\varepsilon<1/2\),
\begin{align}\label{B:eq:per-d}
 &\left|\sum_{\substack{Q\in\cF_d^I\\q\mid n(Q)}}\psi_q(u_Q)
             -g_{c,d}(q)\mathfrak M_d\right|\notag\\
 &\quad\ll_\varepsilon \cL^C q\kappa_c(q)^{1/2}(\#\Lambda_d(1))^{1/2}
   \left\{A\sqrt d\left(1+\frac A{qF}\right)
                     +\frac{F^{1+\varepsilon}}{\sqrt d}\right\}^{1/2}.
\end{align}
In particular the factor \(\kappa_c(q)^{1/2}\) is absent when \((q,c)=1\),
as in the sieve application. There is no lower bound on \(F/A\).
\end{theorem}
\begin{proof}
Corollary~\ref{L:4.4} applies to the sieved subset, because separation
is preserved by passage to a subset and by the dilation to \(u\).
Lemma~\ref{L:6.1} identifies its mean as \(g_{c,d}(q)\mathfrak M_d\).
In Proposition~\ref{L:5.1},
\[
 \frac\lambda Y\asymp A\sqrt d,\qquad
 \frac{\lambda^2}Y\asymp A\sqrt d\frac A{qF},\qquad
 \frac{q^{1/2}}{MY}\lambda^{-\varepsilon}
 \ll\frac{F^{1+\varepsilon}}{\sqrt d}
                 q^{-1/2+\varepsilon}A^{-\varepsilon}
 \le\frac{F^{1+\varepsilon}}{\sqrt d}.
\]
The remaining factor from Sobolev duality is \((\#\Lambda_d(q))^{1/2}\).
For \(\ell\nmid c\),
\[
 |\SL_2(\mathbb F_\ell)|g_{c,d}(\ell)
       =(\ell-\chi_d(\ell))(\ell-1)\le\ell^2;
\]
for \(\ell\mid c\) it is either zero or \(2\ell(\ell-1)\le2\ell^2\).
The exact identity from Lemma~\ref{L:6.1} therefore gives
\(\#\Lambda_d(q)\le\kappa_c(q)q^2\#\Lambda_d(1)\), also for \(q=1\).
Taking square roots proves \eqref{B:eq:per-d}.
\end{proof}

The term \(A/(qF)\) is the cost of wrapping around several periods;
it carries no \(N^\varepsilon\) loss. One may instead use the \(e\)-cusp:
\(w\mapsto-1/(dw)\) sends the point to
\((2a+i/\sqrt d)/e\), and reflection gives the same boxes.
Interchanging \(e,f\) replaces the sieve equation \(A=cB\) by
\(C/d=cB\). The involution
\((A,B,C)\mapsto(C/d,B,dA)\) is a bijection of the quadric modulo every
prime not dividing \(2d\); hence the densities and the estimate are
unchanged.

\begin{lemma}[Class-number bound]\label{L:6.3}
\status{proved}. With \(h(-4d)\) the number of primitive proper classes
of discriminant \(-4d\), and
\(r(d)=\prod_{p^k\parallel d}p^{\lfloor k/2\rfloor}\),
\[
 \#\Lambda_d(1)\le6h(-4d)r(d),\qquad
 \sum_{d\le D}\#\Lambda_d(1)\ll D^{3/2}\log^2(2D).
\]
\end{lemma}
\begin{proof}
Every Type~I form is primitive: a prime dividing \(f\) divides neither
\(2ad\) nor \(4ad\), by \(ef-4a^2d=1\).
For a primitive class represented by \(Q_0\), its admissible level-\(d\)
cosets are bounded by the number of
\((t:s)\in\mathbb P^1(\Z/d\Z)\) with \(Q_0(t,s)\equiv0\pmod d\).
Indeed, the second column of a matrix representing a coset modulo
\(\Gamma^0(d)\) gives this projective point; the condition that the
transformed last coefficient be divisible by \(d\) is necessary.
Passing from \(\Gamma_0(d)\) to \(\Gamma'_1\) costs index at most six
(exactly six for odd \(d\), four for even \(d\)); imposing Type~I parity
can only decrease the count. Stabiliser weights only improve this upper bound.

For \(p^k\parallel d\), make an invertible change of variables so that
\(Q_0=[A_0,B_0,C_0]\) has \(p\nmid A_0\). This is possible by primitivity;
at 2, \(B_0\) is even and one of the outer coefficients is odd.
Points with \(p\mid s\) cannot be zeros. At an odd prime, zeros
\((t:1)\) satisfy
\((2A_0t+B_0)^2\equiv0\pmod{p^k}\), since the discriminant is \(-4d\).
There are exactly \(p^{\lfloor k/2\rfloor}\) possibilities.
At 2, write \(B_0=2B'_0\); the identity
\[
 A_0Q_0(t,1)=(A_0t+B'_0)^2+d
\]
shows that the condition is
\(A_0t+B'_0\equiv0\pmod{2^{\lceil k/2\rceil}}\), again giving
\(2^{\lfloor k/2\rfloor}\) possibilities. CRT proves the first assertion.
The classical imaginary-quadratic class-number formula
\cite{IK} is
\(h(D_0)=w(D_0)\sqrt{|D_0|}L(1,\chi_{D_0})/(2\pi)\),
where \(w(D_0)\le6\) and \(\chi_{D_0}\) is the Kronecker character,
possibly imprimitive. Its nonprincipal character sums are bounded by
\(|D_0|\); splitting the Dirichlet series there and using partial
summation gives \(L(1,\chi_{D_0})\ll\log(2|D_0|)\).
Thus \(h(-4d)\ll\sqrt d\log(2d)\), including nonfundamental discriminants.
Finally,
\[
 \sum_{d\le D}r(d)
 \le\sum_{m\le\sqrt D}m\,\#\{d\le D:m^2\mid d\}
 \le D\sum_{m\le\sqrt D}\frac1m\ll D\log(2D).
\]
Multiplying by \(\sqrt D\log(2D)\) proves the second assertion.
\end{proof}

There is no uniform per-\(d\) relative-error assertion: the orbital
main term contains a possibly small class number. Section~\ref{sec:assembly}
uses only absolute errors and
\(\sum_{d\asymp D}\sqrt{\#\Lambda_d(1)}\ll D^{5/4}\log(2D)\),
by Cauchy--Schwarz and Lemma~\ref{L:6.3}. Relative to the box mass
\(AD\), the first summed error in \eqref{B:eq:per-d} has size
\(\cL^Cq(D/A)^{1/2}(1+A/(qF))^{1/2}\) when \((q,c)=1\).
This is an averaged absolute-error calculation, not a lower bound for
an individual \(\mathfrak M_d\).
